% !PRINTED 121
\weberpage{121}

\begin{center}
{\large \S.\ 38.}

\vspace{0.5em}
{\large\gesperrt{Proof of the Fundamental Theorem.}}
\end{center}

We now proceed to the proof of the \gesperrt{fundamental theorem of algebra}, namely that every equation $f(x)=0$ has at least one root. We first put forward some general propositions.

1. \gesperrt{If $S$ denotes an arbitrary system of real numbers all greater than a certain positive or negative number $C$ (that is, a system which contains no infinitely negative numbers), then there exists a lower bound for the numbers $S$.}

By a \gesperrt{lower bound} $g$ one is to understand a number which is not undercut by any number of the system $S$, but which is of such a nature that, if $\delta$ is an arbitrarily given positive quantity, then between $g$ and $g+\delta$ (including the endpoints) there always lies at least one number of the system $S$, however small $\delta$ may be.

The proof follows directly from the possibility of a cut $(\mathfrak A,\mathfrak B)$, constructed as follows: one throws a number $\alpha$ into $\mathfrak A$ if it is not undercut by any number of the system $S$, and a number $\beta$ into $\mathfrak B$ if it is undercut by at least one number in $S$. The number $g$ determined by this cut is not undercut by any number in $S$, for otherwise there would also be numbers smaller than $g$, and therefore belonging to $\mathfrak A$, which nevertheless are undercut by numbers in $S$; while, on the other hand, every number $\beta$ lying even ever so little above $g$ belongs to $\mathfrak B$ and is therefore undercut by a number in $S$. These are precisely the characteristic marks of the \gesperrt{lower bound}.

2. \gesperrt{If $S'$ is a part of $S$, then the numbers $S'$ also have a lower bound $g'$, and this is either equal to $g$ or greater than $g$.}

For it cannot be smaller than $g$, since otherwise numbers would occur in $S'$, and therefore also in $S$, which lie below $g$.

Now let $f(x)$ be a real function of $x$ which has only finite values in the interval
\begin{equation}
 a\le x\le b
\tag{1}
\end{equation}
%CONT>>
